\documentclass[conference]{IEEEtran}
\usepackage{amsmath,amssymb}
\usepackage{cite}

\begin{document}

\title{Quantum-Classical Architectures for Ethereum Mempool Forecasting}
\author{\IEEEauthorblockN{Quantum Sentinel Vault Team}}
\maketitle

\begin{abstract}
We present a hybrid architecture in which a four-qubit variational quantum circuit
encodes mempool features and feeds a classical classification head to forecast
transaction congestion and adversarial ordering behaviour.
\end{abstract}

\begin{IEEEkeywords}
variational quantum circuits, mempool, Ethereum, hybrid learning
\end{IEEEkeywords}

\section{Introduction}
Public mempools expose pending transactions to ordering attacks~\cite{daian2020flash}.
Variational quantum algorithms~\cite{cerezo2021variational} offer compact models that
can be trained end-to-end with classical networks~\cite{bergholm2018pennylane}.

\section{Architecture}
Classical features $x$ are projected to angles $\theta=\pi\tanh(Wx)$ and embedded by
rotations. Parameterized layers with CNOT entanglement produce expectations
$\langle Z_i\rangle$, which a linear head maps to class logits.

\section{Conclusion}
The hybrid design is a lightweight baseline for quantum-classical mempool forecasting.

\bibliographystyle{IEEEtran}
\bibliography{references}

\end{document}
